\chapter{Surface layer}\label{ch:sfc}

The surface layer scheme computes the exchange coefficients between the ground and the lowest model level:
\begin{itemize}
  \item the friction velocity $u_*$;
  \item the fluxes of momentum, heat and moisture;
  \item the diagnostic 2-m temperature and humidity and 10-m wind.
\end{itemize}
The land-surface model (\cref{ch:lsm}) and the boundary layer or LES use them. The case uses the revised MM5 scheme
of \citet{jimenez2012} (\code{sf_sfclay_physics = 1}). Its core, \src{phys/physics_mmm/sf_sfclayrev.F90}{sf_sfclayrev_run},
is shared with MPAS and wrapped by \code{phys/module_sf_sfclayrev.F}.

\section{Monin--Obukhov similarity}

Near the ground, the mean profiles of wind and potential temperature depend on height $z$ only through the stability
parameter $\zeta = z/L$. The Obukhov length $L$ measures the height above which buoyancy matters as much as shear:
\begin{equation}
  \frac{\kappa z}{u_*}\pd{U}{z} = \phi_m(\zeta), \qquad
  \frac{\kappa z}{\theta_*}\pd{\theta}{z} = \phi_h(\zeta),
\end{equation}
with the von K\'arm\'an constant $\kappa = 0.4$. Integrated from the roughness length $z_0$ to the first model level
$z_a$:
\begin{equation}
  U(z_a) = \frac{u_*}{\kappa}\Bigl[\ln\frac{z_a + z_0}{z_0} - \psi_m\Bigl(\frac{z_a + z_0}{L}\Bigr) + \psi_m\Bigl(\frac{z_0}{L}\Bigr)\Bigr],
\end{equation}
and the same for $\theta$ with $\psi_h$ and the thermal roughness length.

\section{The stability functions}

The revised scheme uses functions that are valid over the full range of stability \citep{jimenez2012}. In the stable
case ($\zeta > 0$), \src{phys/physics_mmm/sf_sfclayrev.F90}{psim_stable_full} and \code{psih_stable_full} are
\begin{equation}
  \psi_m = -6.1\,\ln\!\Bigl(\zeta + \bigl(1 + \zeta^{2.5}\bigr)^{1/2.5}\Bigr), \qquad
  \psi_h = -5.3\,\ln\!\Bigl(\zeta + \bigl(1 + \zeta^{1.1}\bigr)^{1/1.1}\Bigr).
\end{equation}
In the unstable case, the function blends the Kansas-type form, with $x = (1 - 16\zeta)^{1/4}$,
\begin{equation}
  \psi_{m,K} = 2\ln\frac{1+x}{2} + \ln\frac{1+x^2}{2} - 2\arctan x + \frac{\pi}{2},
\end{equation}
with a free-convection form $\psi_{m,C}$ in $y = (1 - 10\zeta)^{0.33}$, weighted by $\zeta^2$:
\begin{equation}
  \psi_m = \frac{\psi_{m,K} + \zeta^2\,\psi_{m,C}}{1 + \zeta^2},
\end{equation}
and similarly for heat (\code{psim_unstable_full}, \code{psih_unstable_full}). For speed, both are tabulated at
initialization on a grid of 1001 values of $\zeta$ (\code{psim_stab}, \code{psih_unstab}, \dots). The run-time
functions \code{psim_stable} etc.\ look values up in these tables inside their range and call the full functions
outside it.

\section{From the bulk Richardson number to \texorpdfstring{$z/L$}{z/L}}

What the model knows at the first level is the bulk Richardson number
\begin{equation}
  \mathrm{Ri}_B = \frac{g}{\theta_v}\,\frac{z_a\,(\theta_{va} - \theta_{vg})}{U_a^2},
\end{equation}
not $L$. The two are related implicitly through the stability functions. The scheme solves for $\zeta$ with the
secant method in \src{phys/physics_mmm/sf_sfclayrev.F90}{zolri}:
\begin{algorithm}[ht]
\caption{\code{zolri}: $\zeta$ from $\mathrm{Ri}_B$ by the secant method.}
\label{alg:zolri}
\begin{algorithmic}[1]
\State $(x_1, x_2) \gets (-5, 0)$ if $\mathrm{Ri}_B < 0$, else $(0, 5)$
\State $f_1 \gets f(x_1)$, $f_2 \gets f(x_2)$, where $f(\zeta) = \mathrm{Ri}(\zeta) - \mathrm{Ri}_B$ (\code{zolri2})
\State result $\gets x_1$
\While{$|x_1 - x_2| > 0.01$ and fewer than 10 iterations}
  \If{$f_1 = f_2$} \textbf{return} the current result \Comment{the case discussed below}
  \EndIf
  \If{$|f_2| < |f_1|$}
    \State $x_1 \gets x_1 - f_1\,(x_2 - x_1)/(f_2 - f_1)$; $f_1 \gets f(x_1)$; result $\gets x_1$
  \Else
    \State $x_2 \gets x_2 - f_2\,(x_2 - x_1)/(f_2 - f_1)$; $f_2 \gets f(x_2)$; result $\gets x_2$
  \EndIf
\EndWhile
\end{algorithmic}
\end{algorithm}

\begin{portnote}
In the original code, if $f_1 = f_2$ on the very first pass, the function returned without assigning its result.
That is undefined behaviour: the value is whatever happens to be in the result register, which can differ between
CPU and GPU. The reproducible reference build defines it: the result is set to the current iterate before the loop
(\cref{alg:zolri}, line 3). A counter (\code{sfclayrev_zolri_undef}) records how often the case occurs; the reference
run of the case records zero (test T-ZOLRI), so the fix changes no result of the case.
\end{portnote}

\section{Fluxes}

With $\zeta$ known, the scheme computes:
\begin{align}
  u_* &= \frac{\kappa\,U_a}{\ln\bigl((z_a + z_0)/z_0\bigr) - \psi_m}, \\
  H &= \rho\,c_p\,C_h\,U_a\,(\theta_g - \theta_a), \qquad
  E = \rho\,C_q\,U_a\,M\,(q_{s}(T_g) - q_a),
\end{align}
where $C_h$ and $C_q$ are the exchange coefficients for heat and moisture and $M$ is the moisture availability. Over
land, the land-surface model replaces the fluxes with its own, using these coefficients (\cref{ch:lsm}). A
convective velocity scale enters $U_a$ in unstable conditions, so the fluxes do not vanish in calm air. The scheme also
interpolates the 2-m and 10-m diagnostics with the same similarity profiles.

\begin{portnote}
The surface layer is a per-point calculation: only the lowest model level is used. On the GPU it is one thread per
$(i,j)$ point (kernel K-SFCLAY).
\begin{itemize}
  \item The tables of $\psi$ are uploaded to the device once.
  \item \code{zolri}, \code{zolri2} and the stability functions are device routines.
  \item Their logarithms, powers and arctangents use the reproducible library, so the iteration takes the same path,
    with the same number of steps, on CPU and GPU.
\end{itemize}
\end{portnote}
